\title{INIZIO DEANGELIS}
\newpage

\chapter{Panoramica del modulo di Software Testing}

\section{Modalità di esame}

Il lavoro richiesto è \textbf{strettamente individuale} e consiste nell'applicare le tecniche del corso a \textbf{due classi differenti} appartenenti a \textbf{un progetto open-source} della \textbf{Apache Software Foundation}.

I sorgenti del progetto devono essere disponibili su \textbf{GitHub} e lo studente deve configurare l'ambiente di lavoro del progetto.

\begin{notaprof}
Si consiglia fortemente di partire da un \emph{fork} personale del progetto su GitHub. In questo modo, se durante il lavoro si individua un bug reale, è anche possibile aprire una \emph{pull request} verso il progetto originale.
\end{notaprof}

\begin{notaesame}
È richiesta la configurazione di un framework di \textbf{Continuous Integration}, ad esempio GitHub Actions o TravisCI. È possibile lavorare solo in locale, ma un approccio di questo tipo avrà un impatto negativo sulla valutazione finale.
\end{notaesame}

\section{Attività richieste: generazione dei test}

\subsection{Scelta delle classi}

\begin{notaesame}
Il docente raccomanda esplicitamente di non scegliere classi banali, né classi legate alla GUI. In questi casi sarebbe difficile dimostrare l'apprendimento dei concetti del corso.
\end{notaesame}

\begin{notaslide}
Se si intuisce che una classe raggiungerà subito il 100\% di coverage senza sforzo, è consigliato contattare il docente.
\end{notaslide}

\subsection{Test manuali}

Su ciascuna classe bisogna definire un insieme di test tramite \textbf{category partition}, adottando preferibilmente un approccio \textbf{black-box} basato sulle funzionalità esplicite o implicite della classe.

\subsection{Test automatici}

Occorre poi generare automaticamente altri test usando tre approcci distinti:
\begin{itemize}
    \item approcci \textbf{randomici};
    \item approcci basati su \textbf{LLM} e sperimentazione di diversi prompt;
    \item approcci guidati da \textbf{coverage control-flow}.
\end{itemize}

\subsection{Integrazione nel build}

\begin{notaprof}
Per evitare problemi di configurazione e interferenze, è consigliabile disabilitare o rimuovere tutti i test nativi del progetto Apache, lasciando nel ciclo di build soltanto i test sviluppati da voi.
\end{notaprof}

\section{Valutazione della qualità dei test}

La qualità dei test deve essere validata scegliendo almeno \textbf{due metriche di adeguatezza}, non necessariamente calcolate da tool automatici, e confrontando i valori ottenuti.

Occorre inoltre:
\begin{itemize}
    \item valutare la significatività dei test rispetto alle funzionalità attese;
    \item migliorare i test manuali aumentando la loro adeguatezza;
    \item applicare tecniche di \textbf{mutation testing};
    \item aggiungere nuovi test per migliorare la robustezza della suite;
    \item stimare la \textbf{reliability} del sistema assumendo un profilo operazionale uniforme.
\end{itemize}

\begin{notaslide}
Per consentire il completamento della build, quando necessario è possibile disabilitare temporaneamente i test che falliscono, purché la cosa venga documentata.
\end{notaslide}

\begin{notaprof}
Se un test fallisce sulla classe originale, ci sono due possibilità: o il test è stato scritto male e quindi va corretto, oppure avete scoperto un bug reale, che va documentato e gestito separatamente.
\end{notaprof}

\section{Varianti delle classi e uso degli LLM}

La parte finale del progetto riguarda l'evoluzione del sistema e il comportamento dei test quando il codice cambia.

A partire dalla classe originale \(C_0\), bisogna chiedere a un LLM di generare quattro varianti di refactoring:
\[
C_1,\; C_2,\; C_3,\; C_4
\]

\begin{notaprof}
Il fulcro dell'esperimento è fornire all'LLM quantità crescenti di informazione. In particolare:
\begin{itemize}
    \item \(C_1\): si fornisce solo la classe originale e si chiede il refactoring;
    \item \(C_2\): si chiede il refactoring imponendo il passaggio dei test black-box;
    \item \(C_3\): si aggiunge anche il vincolo di passare i test di control flow;
    \item \(C_4\): si impone il passaggio di tutti i test, inclusa la mutation analysis.
\end{itemize}
\end{notaprof}

Occorre poi:
\begin{itemize}
    \item eseguire tutti i test originari sulle varianti;
    \item verificare se le funzionalità iniziali siano state mantenute;
    \item generare nuovi test automatici per le varianti;
    \item confrontare qualitativamente i test delle varianti con quelli della classe originale.
\end{itemize}

\begin{notaprof}
Se la variante \(C_4\) non supera un test che era stato fornito esplicitamente come vincolo all'LLM, significa che l'LLM ha fallito il refactoring.
\end{notaprof}

\section{Report finale e consegna}

Lo studente deve redigere un \textbf{report dettagliato} in cui descrive:
\begin{itemize}
    \item le attività svolte;
    \item le decisioni prese;
    \item la metodologia seguita;
    \item i problemi affrontati;
    \item i risultati ottenuti.
\end{itemize}

Il report ideale ha una lunghezza di circa \textbf{12 pagine A4}, a interlinea singola, carattere 10, font Arial, in formato PDF.

\begin{notaprof}
Le 12 pagine si riferiscono al testo descrittivo vero e proprio. Figure, tabelle e listati di codice vanno inseriti alla fine, come allegati, e referenziati nel testo.
\end{notaprof}

\begin{notaesame}
Non bisogna usare figure, tabelle o didascalie come scorciatoia per spostare fuori dal conteggio l'analisi testuale. Le tabelle mostrano i dati; la discussione e l'interpretazione devono stare nel corpo del report.
\end{notaesame}

\subsection{File \texttt{classes.txt}}

Va inoltre consegnato un file di testo chiamato \texttt{classes.txt}, contenente i nomi completi delle classi considerate, in ordine alfabetico, nel formato:
\[
\langle nomePackage \rangle . \langle nomeSubPackage \rangle . \dots . \langle nomeClasse \rangle
\]

\subsection{Invio e discussione orale}

Il report e il file \texttt{classes.txt} vanno inviati via email al docente entro le scadenze pubblicate su Teams. L'esame orale consiste nella presentazione e discussione del report.

\begin{notaprof}
I link al repository GitHub e al framework di CI dovrebbero essere inseriti esplicitamente nel report oppure nella mail di consegna.
\end{notaprof}

\section{Messaggi chiave della lezione}

\begin{notaesame}
Il docente valuterà la comprensione critica del lavoro svolto. Non basta riportare un valore numerico, per esempio un coverage del 37\%: bisogna spiegare perché si è ottenuto quel valore e quali implicazioni ha.
\end{notaesame}

\begin{notaprof}
Quando si generano le varianti \(C_1\)--\(C_4\) con l'LLM, conviene sempre iniziare una nuova sessione di chat fresca, in modo da evitare che il modello venga influenzato dai prompt precedenti e produca varianti troppo simili tra loro.
\end{notaprof}

\subsection*{In sintesi}

Il modulo di Software Testing prevede un esame basato su un progetto individuale rigoroso. Lo studente deve lavorare su due classi non banali di un progetto Apache, generare e valutare test con approcci diversi, integrare i test in una pipeline di build, studiare l'effetto delle mutazioni e delle varianti generate con LLM, e infine documentare tutto in un report tecnico dettagliato da discutere oralmente.